% Options for packages loaded elsewhere
% Options for packages loaded elsewhere
\PassOptionsToPackage{unicode}{hyperref}
\PassOptionsToPackage{hyphens}{url}
\PassOptionsToPackage{dvipsnames,svgnames,x11names}{xcolor}
%
\documentclass[
  letterpaper,
  DIV=11,
  numbers=noendperiod]{scrartcl}
\usepackage{xcolor}
\usepackage{amsmath,amssymb}
\setcounter{secnumdepth}{5}
\usepackage{iftex}
\ifPDFTeX
  \usepackage[T1]{fontenc}
  \usepackage[utf8]{inputenc}
  \usepackage{textcomp} % provide euro and other symbols
\else % if luatex or xetex
  \usepackage{unicode-math} % this also loads fontspec
  \defaultfontfeatures{Scale=MatchLowercase}
  \defaultfontfeatures[\rmfamily]{Ligatures=TeX,Scale=1}
\fi
\usepackage{lmodern}
\ifPDFTeX\else
  % xetex/luatex font selection
\fi
% Use upquote if available, for straight quotes in verbatim environments
\IfFileExists{upquote.sty}{\usepackage{upquote}}{}
\IfFileExists{microtype.sty}{% use microtype if available
  \usepackage[]{microtype}
  \UseMicrotypeSet[protrusion]{basicmath} % disable protrusion for tt fonts
}{}
\makeatletter
\@ifundefined{KOMAClassName}{% if non-KOMA class
  \IfFileExists{parskip.sty}{%
    \usepackage{parskip}
  }{% else
    \setlength{\parindent}{0pt}
    \setlength{\parskip}{6pt plus 2pt minus 1pt}}
}{% if KOMA class
  \KOMAoptions{parskip=half}}
\makeatother
% Make \paragraph and \subparagraph free-standing
\makeatletter
\ifx\paragraph\undefined\else
  \let\oldparagraph\paragraph
  \renewcommand{\paragraph}{
    \@ifstar
      \xxxParagraphStar
      \xxxParagraphNoStar
  }
  \newcommand{\xxxParagraphStar}[1]{\oldparagraph*{#1}\mbox{}}
  \newcommand{\xxxParagraphNoStar}[1]{\oldparagraph{#1}\mbox{}}
\fi
\ifx\subparagraph\undefined\else
  \let\oldsubparagraph\subparagraph
  \renewcommand{\subparagraph}{
    \@ifstar
      \xxxSubParagraphStar
      \xxxSubParagraphNoStar
  }
  \newcommand{\xxxSubParagraphStar}[1]{\oldsubparagraph*{#1}\mbox{}}
  \newcommand{\xxxSubParagraphNoStar}[1]{\oldsubparagraph{#1}\mbox{}}
\fi
\makeatother

\usepackage{color}
\usepackage{fancyvrb}
\newcommand{\VerbBar}{|}
\newcommand{\VERB}{\Verb[commandchars=\\\{\}]}
\DefineVerbatimEnvironment{Highlighting}{Verbatim}{commandchars=\\\{\}}
% Add ',fontsize=\small' for more characters per line
\usepackage{framed}
\definecolor{shadecolor}{RGB}{241,243,245}
\newenvironment{Shaded}{\begin{snugshade}}{\end{snugshade}}
\newcommand{\AlertTok}[1]{\textcolor[rgb]{0.68,0.00,0.00}{#1}}
\newcommand{\AnnotationTok}[1]{\textcolor[rgb]{0.37,0.37,0.37}{#1}}
\newcommand{\AttributeTok}[1]{\textcolor[rgb]{0.40,0.45,0.13}{#1}}
\newcommand{\BaseNTok}[1]{\textcolor[rgb]{0.68,0.00,0.00}{#1}}
\newcommand{\BuiltInTok}[1]{\textcolor[rgb]{0.00,0.23,0.31}{#1}}
\newcommand{\CharTok}[1]{\textcolor[rgb]{0.13,0.47,0.30}{#1}}
\newcommand{\CommentTok}[1]{\textcolor[rgb]{0.37,0.37,0.37}{#1}}
\newcommand{\CommentVarTok}[1]{\textcolor[rgb]{0.37,0.37,0.37}{\textit{#1}}}
\newcommand{\ConstantTok}[1]{\textcolor[rgb]{0.56,0.35,0.01}{#1}}
\newcommand{\ControlFlowTok}[1]{\textcolor[rgb]{0.00,0.23,0.31}{\textbf{#1}}}
\newcommand{\DataTypeTok}[1]{\textcolor[rgb]{0.68,0.00,0.00}{#1}}
\newcommand{\DecValTok}[1]{\textcolor[rgb]{0.68,0.00,0.00}{#1}}
\newcommand{\DocumentationTok}[1]{\textcolor[rgb]{0.37,0.37,0.37}{\textit{#1}}}
\newcommand{\ErrorTok}[1]{\textcolor[rgb]{0.68,0.00,0.00}{#1}}
\newcommand{\ExtensionTok}[1]{\textcolor[rgb]{0.00,0.23,0.31}{#1}}
\newcommand{\FloatTok}[1]{\textcolor[rgb]{0.68,0.00,0.00}{#1}}
\newcommand{\FunctionTok}[1]{\textcolor[rgb]{0.28,0.35,0.67}{#1}}
\newcommand{\ImportTok}[1]{\textcolor[rgb]{0.00,0.46,0.62}{#1}}
\newcommand{\InformationTok}[1]{\textcolor[rgb]{0.37,0.37,0.37}{#1}}
\newcommand{\KeywordTok}[1]{\textcolor[rgb]{0.00,0.23,0.31}{\textbf{#1}}}
\newcommand{\NormalTok}[1]{\textcolor[rgb]{0.00,0.23,0.31}{#1}}
\newcommand{\OperatorTok}[1]{\textcolor[rgb]{0.37,0.37,0.37}{#1}}
\newcommand{\OtherTok}[1]{\textcolor[rgb]{0.00,0.23,0.31}{#1}}
\newcommand{\PreprocessorTok}[1]{\textcolor[rgb]{0.68,0.00,0.00}{#1}}
\newcommand{\RegionMarkerTok}[1]{\textcolor[rgb]{0.00,0.23,0.31}{#1}}
\newcommand{\SpecialCharTok}[1]{\textcolor[rgb]{0.37,0.37,0.37}{#1}}
\newcommand{\SpecialStringTok}[1]{\textcolor[rgb]{0.13,0.47,0.30}{#1}}
\newcommand{\StringTok}[1]{\textcolor[rgb]{0.13,0.47,0.30}{#1}}
\newcommand{\VariableTok}[1]{\textcolor[rgb]{0.07,0.07,0.07}{#1}}
\newcommand{\VerbatimStringTok}[1]{\textcolor[rgb]{0.13,0.47,0.30}{#1}}
\newcommand{\WarningTok}[1]{\textcolor[rgb]{0.37,0.37,0.37}{\textit{#1}}}

\usepackage{longtable,booktabs,array}
\usepackage{calc} % for calculating minipage widths
% Correct order of tables after \paragraph or \subparagraph
\usepackage{etoolbox}
\makeatletter
\patchcmd\longtable{\par}{\if@noskipsec\mbox{}\fi\par}{}{}
\makeatother
% Allow footnotes in longtable head/foot
\IfFileExists{footnotehyper.sty}{\usepackage{footnotehyper}}{\usepackage{footnote}}
\makesavenoteenv{longtable}
\usepackage{graphicx}
\makeatletter
\newsavebox\pandoc@box
\newcommand*\pandocbounded[1]{% scales image to fit in text height/width
  \sbox\pandoc@box{#1}%
  \Gscale@div\@tempa{\textheight}{\dimexpr\ht\pandoc@box+\dp\pandoc@box\relax}%
  \Gscale@div\@tempb{\linewidth}{\wd\pandoc@box}%
  \ifdim\@tempb\p@<\@tempa\p@\let\@tempa\@tempb\fi% select the smaller of both
  \ifdim\@tempa\p@<\p@\scalebox{\@tempa}{\usebox\pandoc@box}%
  \else\usebox{\pandoc@box}%
  \fi%
}
% Set default figure placement to htbp
\def\fps@figure{htbp}
\makeatother


% definitions for citeproc citations
\NewDocumentCommand\citeproctext{}{}
\NewDocumentCommand\citeproc{mm}{%
  \begingroup\def\citeproctext{#2}\cite{#1}\endgroup}
\makeatletter
 % allow citations to break across lines
 \let\@cite@ofmt\@firstofone
 % avoid brackets around text for \cite:
 \def\@biblabel#1{}
 \def\@cite#1#2{{#1\if@tempswa , #2\fi}}
\makeatother
\newlength{\cslhangindent}
\setlength{\cslhangindent}{1.5em}
\newlength{\csllabelwidth}
\setlength{\csllabelwidth}{3em}
\newenvironment{CSLReferences}[2] % #1 hanging-indent, #2 entry-spacing
 {\begin{list}{}{%
  \setlength{\itemindent}{0pt}
  \setlength{\leftmargin}{0pt}
  \setlength{\parsep}{0pt}
  % turn on hanging indent if param 1 is 1
  \ifodd #1
   \setlength{\leftmargin}{\cslhangindent}
   \setlength{\itemindent}{-1\cslhangindent}
  \fi
  % set entry spacing
  \setlength{\itemsep}{#2\baselineskip}}}
 {\end{list}}
\usepackage{calc}
\newcommand{\CSLBlock}[1]{\hfill\break\parbox[t]{\linewidth}{\strut\ignorespaces#1\strut}}
\newcommand{\CSLLeftMargin}[1]{\parbox[t]{\csllabelwidth}{\strut#1\strut}}
\newcommand{\CSLRightInline}[1]{\parbox[t]{\linewidth - \csllabelwidth}{\strut#1\strut}}
\newcommand{\CSLIndent}[1]{\hspace{\cslhangindent}#1}



\setlength{\emergencystretch}{3em} % prevent overfull lines

\providecommand{\tightlist}{%
  \setlength{\itemsep}{0pt}\setlength{\parskip}{0pt}}



 


\renewcommand{\thefigure}{S\arabic{figure}}
\renewcommand{\thesection}{S\arabic{section}}
\KOMAoption{captions}{tableheading}
\makeatletter
\@ifpackageloaded{caption}{}{\usepackage{caption}}
\AtBeginDocument{%
\ifdefined\contentsname
  \renewcommand*\contentsname{Table of contents}
\else
  \newcommand\contentsname{Table of contents}
\fi
\ifdefined\listfigurename
  \renewcommand*\listfigurename{List of Figures}
\else
  \newcommand\listfigurename{List of Figures}
\fi
\ifdefined\listtablename
  \renewcommand*\listtablename{List of Tables}
\else
  \newcommand\listtablename{List of Tables}
\fi
\ifdefined\figurename
  \renewcommand*\figurename{Supp. Figure}
\else
  \newcommand\figurename{Supp. Figure}
\fi
\ifdefined\tablename
  \renewcommand*\tablename{Table}
\else
  \newcommand\tablename{Table}
\fi
}
\@ifpackageloaded{float}{}{\usepackage{float}}
\floatstyle{ruled}
\@ifundefined{c@chapter}{\newfloat{codelisting}{h}{lop}}{\newfloat{codelisting}{h}{lop}[chapter]}
\floatname{codelisting}{Listing}
\newcommand*\listoflistings{\listof{codelisting}{List of Listings}}
\makeatother
\makeatletter
\makeatother
\makeatletter
\@ifpackageloaded{caption}{}{\usepackage{caption}}
\@ifpackageloaded{subcaption}{}{\usepackage{subcaption}}
\makeatother
\usepackage{bookmark}
\IfFileExists{xurl.sty}{\usepackage{xurl}}{} % add URL line breaks if available
\urlstyle{same}
\hypersetup{
  pdftitle={Supplament to ``Spatial aggregation confounds null model-based inference of species association networks''},
  pdfauthor={Andrew J. Rominger},
  colorlinks=true,
  linkcolor={blue},
  filecolor={Maroon},
  citecolor={Blue},
  urlcolor={Blue},
  pdfcreator={LaTeX via pandoc}}


\title{Supplament to ``Spatial aggregation confounds null model-based
inference of species association networks''}
\author{Andrew J. Rominger}
\date{}
\begin{document}
\maketitle


\section{Reproducibility}\label{reproducibility}

The results of this study can be fully reproduced by installing the
package \emph{SpatAggNet} written in R (R Core Team 2018) and available
on GitHub. \emph{SpatAggNet} has a dependency on the custom pacakge
\emph{pika} also available from GitHub. \emph{pika} must be installed
first for \emph{SpatAggNet} to install.

\begin{Shaded}
\begin{Highlighting}[]
\NormalTok{pika\_pack }\OtherTok{\textless{}{-}} \FunctionTok{require}\NormalTok{(pika)}
\ControlFlowTok{if}\NormalTok{(}\SpecialCharTok{!}\NormalTok{pika\_pack) }\FunctionTok{install\_github}\NormalTok{(}\StringTok{"ajrominger/pika"}\NormalTok{)}

\NormalTok{this\_pack }\OtherTok{\textless{}{-}} \FunctionTok{require}\NormalTok{(SpatAggNet)}
\ControlFlowTok{if}\NormalTok{(}\SpecialCharTok{!}\NormalTok{this\_pack) }\FunctionTok{install\_github}\NormalTok{(}\StringTok{"ajrominger/SpatAggNet"}\NormalTok{)}
\end{Highlighting}
\end{Shaded}

The \emph{SpatAggNet} package includes documented (Wickham \emph{et al.}
2018) and unit-tested (Wickham 2011) functions to carry out all
analyses.

The following R computing environment is used for all simulations and
analyses:

\begin{Shaded}
\begin{Highlighting}[]
\FunctionTok{library}\NormalTok{(SpatAggNet)}
\FunctionTok{library}\NormalTok{(pika)}
\FunctionTok{library}\NormalTok{(dplyr)}
\FunctionTok{library}\NormalTok{(tidyr)}
\FunctionTok{library}\NormalTok{(ggplot2)}
\FunctionTok{library}\NormalTok{(cowplot)}
\FunctionTok{library}\NormalTok{(knitr)}
\FunctionTok{library}\NormalTok{(parallel)}

\CommentTok{\# threading defaults}
\NormalTok{nthrd }\OtherTok{\textless{}{-}} \FunctionTok{detectCores}\NormalTok{()}
\NormalTok{nthrd }\OtherTok{\textless{}{-}} \FunctionTok{ifelse}\NormalTok{(}\FunctionTok{round}\NormalTok{(nthrd }\SpecialCharTok{*} \FloatTok{0.8}\NormalTok{) }\SpecialCharTok{\textgreater{}=}\NormalTok{ nthrd }\SpecialCharTok{{-}} \DecValTok{1} \SpecialCharTok{\&}\NormalTok{ nthrd }\SpecialCharTok{\textgreater{}} \DecValTok{1}\NormalTok{, }
\NormalTok{                nthrd }\SpecialCharTok{{-}} \DecValTok{1}\NormalTok{, }
                \FunctionTok{round}\NormalTok{(nthrd }\SpecialCharTok{*} \FloatTok{0.8}\NormalTok{))}


\FunctionTok{sessionInfo}\NormalTok{() }\SpecialCharTok{|\textgreater{}} 
    \FunctionTok{print}\NormalTok{(}\AttributeTok{locale =} \ConstantTok{FALSE}\NormalTok{)}
\end{Highlighting}
\end{Shaded}

\begin{verbatim}
R version 4.5.2 (2025-10-31)
Platform: aarch64-apple-darwin20
Running under: macOS Sequoia 15.0.1

Matrix products: default
BLAS:   /System/Library/Frameworks/Accelerate.framework/Versions/A/Frameworks/vecLib.framework/Versions/A/libBLAS.dylib 
LAPACK: /Library/Frameworks/R.framework/Versions/4.5-arm64/Resources/lib/libRlapack.dylib;  LAPACK version 3.12.1

attached base packages:
[1] parallel  stats     graphics  grDevices utils     datasets  methods  
[8] base     

other attached packages:
[1] knitr_1.50         cowplot_1.2.0.9000 ggplot2_4.0.1      tidyr_1.3.1       
[5] dplyr_1.1.4        SpatAggNet_0.1.0   pika_1.0          

loaded via a namespace (and not attached):
 [1] vctrs_0.6.5        cli_3.6.5          rlang_1.1.6        xfun_0.52         
 [5] purrr_1.1.0        generics_0.1.4     S7_0.2.1           jsonlite_2.0.0    
 [9] glue_1.8.0         htmltools_0.5.8.1  scales_1.4.0       rmarkdown_2.29    
[13] grid_4.5.2         evaluate_1.0.4     tibble_3.3.0       fastmap_1.2.0     
[17] yaml_2.3.10        lifecycle_1.0.4    compiler_4.5.2     RColorBrewer_1.1-3
[21] pkgconfig_2.0.3    rstudioapi_0.17.1  farver_2.1.2       digest_0.6.39     
[25] R6_2.6.1           tidyselect_1.2.1   pillar_1.11.0      magrittr_2.0.4    
[29] withr_3.0.2        gtable_0.3.6       tools_4.5.2       
\end{verbatim}

The quarto (Allaire \emph{et al.} 2025) document used to generate this
supplement can be found with the R command
\texttt{system.file("ms/SpatAggNet\_supp.qmd",\ package\ =\ "SpatAggNet")}.

\section{Computing spatial abundance
distributions}\label{computing-spatial-abundance-distributions}

The data included with the \emph{SpatAggNet} package are from Calatayud
\emph{et al.} (2019) and take the form of a list of site by species
matrices.

\begin{Shaded}
\begin{Highlighting}[]
\CommentTok{\# load list of site by species matrices provided in SpatAggNet}
\FunctionTok{data}\NormalTok{(}\StringTok{"abund\_mats"}\NormalTok{)}

\CommentTok{\# remove decimals and empty rows/columns}
\NormalTok{obs\_dat }\OtherTok{\textless{}{-}} \FunctionTok{lapply}\NormalTok{(abund\_mats, }\ControlFlowTok{function}\NormalTok{(x) \{}
\NormalTok{    x }\OtherTok{\textless{}{-}} \FunctionTok{ceiling}\NormalTok{(x)}
\NormalTok{    x }\OtherTok{\textless{}{-}}\NormalTok{ x[}\FunctionTok{rowSums}\NormalTok{(x) }\SpecialCharTok{\textgreater{}} \DecValTok{0}\NormalTok{, }\FunctionTok{colSums}\NormalTok{(x) }\SpecialCharTok{\textgreater{}} \DecValTok{0}\NormalTok{]}
    
    \FunctionTok{return}\NormalTok{(x)}
\NormalTok{\})}
\end{Highlighting}
\end{Shaded}

In order to investigate the affect of intra-specific spatial aggregation
(ISSA) on inferred networks, we first calculate the parameters of the
species abundance distribution (SAD) for each data set and the
intra-specific spatial abundance distributions (SSAD) for each species
in each data set. The \texttt{nb\_fit} function from \emph{SpatAggNet}
computes SSAD parameters as well as summary statistics including
goodness of fit.

\begin{Shaded}
\begin{Highlighting}[]
\CommentTok{\# SSAD statistics for each species in each data set}
\NormalTok{ssad\_stats }\OtherTok{\textless{}{-}} \FunctionTok{mclapply}\NormalTok{(obs\_dat, }
                       \AttributeTok{mc.cores =}\NormalTok{ nthrd, }
                       \AttributeTok{FUN =} \ControlFlowTok{function}\NormalTok{(x) \{}
\NormalTok{                           nb\_info }\OtherTok{\textless{}{-}} \FunctionTok{nb\_fit}\NormalTok{(x)}
                           \FunctionTok{cbind}\NormalTok{(}\AttributeTok{nsite =} \FunctionTok{nrow}\NormalTok{(x), }\AttributeTok{nspp =} \FunctionTok{ncol}\NormalTok{(x), }
                                 \AttributeTok{J =} \FunctionTok{sum}\NormalTok{(x), nb\_info)}
\NormalTok{                       \}) }\SpecialCharTok{|\textgreater{}} 
    \FunctionTok{bind\_rows}\NormalTok{(}\AttributeTok{.id =} \StringTok{"study"}\NormalTok{)}
\end{Highlighting}
\end{Shaded}

\subsection{Prevelance of the negative binomial
SSAD}\label{prevelance-of-the-negative-binomial-ssad}

We would like to know how prevalent the negative binomial SSAD is in the
data studied here. This is complicated by the fact that the negative
binomial likelihood function does not have a finite maximum when the
mean of the data is greater than or equal to the variance. In the data
studied here, we see the greatest proportion of cases where
\(\mathbb{E}[x] \geq \mathbb{Var}[x]\) for rare species and when the
number of spatial replicates is small Supp. Figure~\ref{fig-inf-k}.
Unstable estimation for rare species with small sample size is to be
expected, but combined with the fact that most species are rare, it is a
challenge to assess the prevalence of the negative binomial SSAD.

\begin{Shaded}
\begin{Highlighting}[]
\NormalTok{abund\_nsite }\OtherTok{\textless{}{-}} \FunctionTok{expand.grid}\NormalTok{(}\AttributeTok{abund =} \DecValTok{1}\SpecialCharTok{:}\DecValTok{40}\NormalTok{, }\AttributeTok{nsite =} \DecValTok{10}\SpecialCharTok{:}\DecValTok{90}\NormalTok{)}

\NormalTok{bad\_k }\OtherTok{\textless{}{-}} \FunctionTok{sapply}\NormalTok{(}\DecValTok{1}\SpecialCharTok{:}\FunctionTok{nrow}\NormalTok{(abund\_nsite), }\ControlFlowTok{function}\NormalTok{(i) \{}
\NormalTok{    j }\OtherTok{\textless{}{-}}\NormalTok{ ssad\_stats}\SpecialCharTok{$}\NormalTok{abund }\SpecialCharTok{\textgreater{}}\NormalTok{ abund\_nsite}\SpecialCharTok{$}\NormalTok{abund[i] }\SpecialCharTok{\&}
\NormalTok{        ssad\_stats}\SpecialCharTok{$}\NormalTok{nsite }\SpecialCharTok{\textgreater{}}\NormalTok{ abund\_nsite}\SpecialCharTok{$}\NormalTok{nsite[i]}
    
    \FunctionTok{mean}\NormalTok{(}\SpecialCharTok{!}\FunctionTok{is.finite}\NormalTok{(ssad\_stats}\SpecialCharTok{$}\NormalTok{size[j]))}
\NormalTok{\})}

\NormalTok{prop\_finite\_k }\OtherTok{\textless{}{-}} \FunctionTok{data.frame}\NormalTok{(abund\_nsite, bad\_k)}


\FunctionTok{ggplot}\NormalTok{(prop\_finite\_k, }\FunctionTok{aes}\NormalTok{(abund, nsite, }\AttributeTok{fill =}\NormalTok{ bad\_k)) }\SpecialCharTok{+}
    \FunctionTok{geom\_tile}\NormalTok{() }\SpecialCharTok{+}
    \FunctionTok{scale\_fill\_viridis\_c}\NormalTok{(}\AttributeTok{name =} \StringTok{"Cumulative}\SpecialCharTok{\textbackslash{}n}\StringTok{prob. of}\SpecialCharTok{\textbackslash{}n}\StringTok{undefined}\SpecialCharTok{\textbackslash{}n}\StringTok{log like"}\NormalTok{) }\SpecialCharTok{+}
    \FunctionTok{theme\_cowplot}\NormalTok{() }\SpecialCharTok{+}
    \FunctionTok{scale\_y\_continuous}\NormalTok{(}\AttributeTok{breaks =} \FunctionTok{c}\NormalTok{(}\DecValTok{10}\NormalTok{, }\DecValTok{30}\NormalTok{, }\DecValTok{50}\NormalTok{, }\DecValTok{70}\NormalTok{, }\DecValTok{90}\NormalTok{)) }\SpecialCharTok{+}
    \FunctionTok{xlab}\NormalTok{(}\StringTok{"Total species abundance"}\NormalTok{) }\SpecialCharTok{+}
    \FunctionTok{ylab}\NormalTok{(}\StringTok{"Number of sites in study"}\NormalTok{)}
\end{Highlighting}
\end{Shaded}

\begin{figure}[H]

\centering{

\pandocbounded{\includegraphics[keepaspectratio]{SpatAggNet_supp_files/figure-pdf/fig-inf-k-1.pdf}}

}

\caption{\label{fig-inf-k}Cumulative probability of encountering a
species with an undefined log likelihood function. Probability
accumulates from right to left and top to bottom. This highlights the
difficulty of fitting a negative binomial distribution to rare species
in data sets with few sites.}

\end{figure}%

I therefore take two approaches to assessing the prevalence of the
negative binomial SSAD:

\begin{enumerate}
\def\labelenumi{\arabic{enumi}.}
\tightlist
\item
  where the likelihood is defined, I compare the support for a negative
  binomial versus a Poisson model with AIC
\item
  for all the data, whether the likelihood is undefined or defined, I
  evaluate the goodness of fit of the negative binomial model to the
  data (described in more detail below)
\end{enumerate}

\begin{Shaded}
\begin{Highlighting}[]
\NormalTok{mod\_select }\OtherTok{\textless{}{-}} \FunctionTok{filter}\NormalTok{(ssad\_stats, }
                     \CommentTok{\# only cases with defined likelihood}
                     \FunctionTok{is.finite}\NormalTok{(size)) }\SpecialCharTok{|\textgreater{}} 
    \CommentTok{\# bin by abundance}
    \FunctionTok{mutate}\NormalTok{(}\AttributeTok{abund\_bin =} \FunctionTok{cut}\NormalTok{(abund, }\FunctionTok{c}\NormalTok{(}\FunctionTok{seq}\NormalTok{(}\DecValTok{1}\NormalTok{, }\DecValTok{800}\NormalTok{, }\AttributeTok{length.out =} \DecValTok{9}\NormalTok{), }
                                    \FunctionTok{exp}\NormalTok{(}\FunctionTok{seq}\NormalTok{(}\FunctionTok{log}\NormalTok{(}\DecValTok{1000}\NormalTok{), }\FunctionTok{log}\NormalTok{(}\FunctionTok{max}\NormalTok{(abund)), }
                                        \AttributeTok{length.out =} \DecValTok{4}\NormalTok{))))) }\SpecialCharTok{|\textgreater{}} 
    \FunctionTok{group\_by}\NormalTok{(abund\_bin) }\SpecialCharTok{|\textgreater{}} 
    \CommentTok{\# calculations to plot bars with nb on top of pois}
    \FunctionTok{summarize}\NormalTok{(}\AttributeTok{pois\_max =} \FunctionTok{mean}\NormalTok{(delta\_AIC }\SpecialCharTok{\textgreater{}=} \DecValTok{0}\NormalTok{),}
              \AttributeTok{nb\_max =} \FunctionTok{mean}\NormalTok{(delta\_AIC }\SpecialCharTok{\textless{}} \DecValTok{0}\NormalTok{) }\SpecialCharTok{+}\NormalTok{ pois\_max, }
              \AttributeTok{pois\_min =} \DecValTok{0}\NormalTok{, }
              \AttributeTok{nb\_min =}\NormalTok{ pois\_max) }\SpecialCharTok{|\textgreater{}} 
    \FunctionTok{ungroup}\NormalTok{() }\SpecialCharTok{|\textgreater{}} 
    \CommentTok{\# resolve upper and lower bounds of rects from}
    \CommentTok{\# output of \textasciigrave{}cut\textasciigrave{}}
    \FunctionTok{pivot\_longer}\NormalTok{(}\AttributeTok{cols =} \SpecialCharTok{{-}}\NormalTok{abund\_bin,}
                 \AttributeTok{names\_to =} \FunctionTok{c}\NormalTok{(}\StringTok{"model"}\NormalTok{, }\StringTok{".value"}\NormalTok{),}
                 \AttributeTok{names\_sep =} \StringTok{"\_"}\NormalTok{) }\SpecialCharTok{|\textgreater{}} 
    \FunctionTok{mutate}\NormalTok{(}\AttributeTok{abund\_max =} \FunctionTok{gsub}\NormalTok{(}\StringTok{".*,|}\SpecialCharTok{\textbackslash{}\textbackslash{}}\StringTok{]"}\NormalTok{, }\StringTok{""}\NormalTok{, abund\_bin) }\SpecialCharTok{|\textgreater{}}
               \FunctionTok{as.numeric}\NormalTok{(),}
           \AttributeTok{abund\_min =} \FunctionTok{gsub}\NormalTok{(}\StringTok{"}\SpecialCharTok{\textbackslash{}\textbackslash{}}\StringTok{(|,.*"}\NormalTok{, }\StringTok{""}\NormalTok{, abund\_bin) }\SpecialCharTok{|\textgreater{}}
               \FunctionTok{as.numeric}\NormalTok{())}

\FunctionTok{ggplot}\NormalTok{(mod\_select, }\FunctionTok{aes}\NormalTok{(}\AttributeTok{xmin =}\NormalTok{ abund\_min, }\AttributeTok{xmax =}\NormalTok{ abund\_max, }
                       \AttributeTok{ymin =}\NormalTok{ min, }\AttributeTok{ymax =}\NormalTok{ max, }\AttributeTok{fill =}\NormalTok{ model)) }\SpecialCharTok{+}
    \FunctionTok{geom\_rect}\NormalTok{(}\AttributeTok{color =} \StringTok{"black"}\NormalTok{, }\AttributeTok{linewidth =} \FloatTok{0.5}\NormalTok{) }\SpecialCharTok{+}
    \FunctionTok{scale\_fill\_discrete}\NormalTok{(}\AttributeTok{name =} \StringTok{"Model favored by AIC"}\NormalTok{, }
                        \AttributeTok{labels =} \FunctionTok{c}\NormalTok{(}\StringTok{"Negative binomial"}\NormalTok{, }\StringTok{"Poisson"}\NormalTok{),}
                        \AttributeTok{palette =} \FunctionTok{hsv}\NormalTok{(}\FunctionTok{c}\NormalTok{(}\FloatTok{0.52}\NormalTok{, }\FloatTok{0.08}\NormalTok{), }\FunctionTok{c}\NormalTok{(}\DecValTok{1}\NormalTok{, }\FloatTok{0.8}\NormalTok{), }\FunctionTok{c}\NormalTok{(}\FloatTok{0.7}\NormalTok{, }\FloatTok{0.9}\NormalTok{))) }\SpecialCharTok{+}
    \FunctionTok{xlab}\NormalTok{(}\StringTok{"Total abundance"}\NormalTok{) }\SpecialCharTok{+}
    \FunctionTok{ylab}\NormalTok{(}\StringTok{"Proportion"}\NormalTok{) }\SpecialCharTok{+}
    \FunctionTok{theme\_cowplot}\NormalTok{() }\SpecialCharTok{+}
    \FunctionTok{theme}\NormalTok{(}\AttributeTok{legend.position =} \FunctionTok{c}\NormalTok{(}\FloatTok{0.92}\NormalTok{, }\FloatTok{0.92}\NormalTok{),}
      \AttributeTok{legend.justification =} \FunctionTok{c}\NormalTok{(}\StringTok{"right"}\NormalTok{, }\StringTok{"top"}\NormalTok{), }
      \AttributeTok{legend.background =} \FunctionTok{element\_rect}\NormalTok{(}\AttributeTok{fill =} \StringTok{"white"}\NormalTok{, }\AttributeTok{color =} \StringTok{"black"}\NormalTok{), }
      \AttributeTok{legend.margin =} \FunctionTok{margin}\NormalTok{(}\AttributeTok{t =} \DecValTok{2}\NormalTok{, }\AttributeTok{r =} \DecValTok{2}\NormalTok{, }\AttributeTok{b =} \DecValTok{2}\NormalTok{, }\AttributeTok{l =} \DecValTok{2}\NormalTok{))}
\end{Highlighting}
\end{Shaded}

\begin{figure}[H]

\centering{

\pandocbounded{\includegraphics[keepaspectratio]{SpatAggNet_supp_files/figure-pdf/fig-aic-ssad-1.pdf}}

}

\caption{\label{fig-aic-ssad}Proportion of species favoring a negative
binomial or Poisson distribution. Model selection is shown across bins
of total species abundance. Bins are 100 wide from 0 to 800, then
increase in width exponentially because fewer species have very large
abundances.}

\end{figure}%

Supp. Figure~\ref{fig-aic-ssad} shows that the negative binomial model
is preferred over the Poisson an overwhelmingly majority of the time.
Only for rare species is the Poisson occasionally favored by AIC. Even
in those cases where Poisson is favored the average
\(\Delta \text{AIC}\) separating Poisson and negative binomial is 1.92
with a maximum of 2.07, meaning even when the Poisson is favored, it is
only marginally favored.

In Supp. Figure~\ref{fig-k-mod} I also investigate the relationship of
the clustering parameter \(k\) with species relative abundance. The
purpose of this analysis is to be able to simulate random but realistic
data.

\begin{Shaded}
\begin{Highlighting}[]
\NormalTok{k\_mod }\OtherTok{\textless{}{-}} \FunctionTok{lm}\NormalTok{(logk }\SpecialCharTok{\textasciitilde{}}\NormalTok{ loga, }
            \AttributeTok{data =} \FunctionTok{filter}\NormalTok{(ssad\_stats, }\FunctionTok{is.finite}\NormalTok{(size)) }\SpecialCharTok{|\textgreater{}} 
                \FunctionTok{mutate}\NormalTok{(}\AttributeTok{loga =} \FunctionTok{log}\NormalTok{(abund }\SpecialCharTok{/}\NormalTok{ J), }
                       \AttributeTok{logk =} \FunctionTok{log}\NormalTok{(size)))}

\CommentTok{\# a function to randomly draw a k value given abundance}
\NormalTok{k\_fun }\OtherTok{\textless{}{-}} \ControlFlowTok{function}\NormalTok{(abund) \{}
\NormalTok{    J }\OtherTok{\textless{}{-}} \FunctionTok{sum}\NormalTok{(abund)}
    
    \FunctionTok{exp}\NormalTok{(}\FunctionTok{predict}\NormalTok{(k\_mod, }\AttributeTok{newdata =} \FunctionTok{data.frame}\NormalTok{(}\AttributeTok{loga =} \FunctionTok{log}\NormalTok{(abund }\SpecialCharTok{/}\NormalTok{ J))) }\SpecialCharTok{+}
            \FunctionTok{rnorm}\NormalTok{(}\FunctionTok{length}\NormalTok{(abund), }\AttributeTok{sd =} \FunctionTok{summary}\NormalTok{(k\_mod)}\SpecialCharTok{$}\NormalTok{sigma))}
\NormalTok{\}}
\end{Highlighting}
\end{Shaded}

\begin{Shaded}
\begin{Highlighting}[]
\FunctionTok{ggplot}\NormalTok{(}\FunctionTok{filter}\NormalTok{(ssad\_stats, }\FunctionTok{is.finite}\NormalTok{(size)), }
       \FunctionTok{aes}\NormalTok{(abund }\SpecialCharTok{/}\NormalTok{ J, size)) }\SpecialCharTok{+}
    \FunctionTok{geom\_point}\NormalTok{(}\AttributeTok{alpha =} \FloatTok{0.25}\NormalTok{) }\SpecialCharTok{+}
    \FunctionTok{scale\_x\_log10}\NormalTok{() }\SpecialCharTok{+}
    \FunctionTok{scale\_y\_log10}\NormalTok{() }\SpecialCharTok{+}
    \FunctionTok{geom\_smooth}\NormalTok{(}\AttributeTok{method =} \StringTok{"lm"}\NormalTok{, }\AttributeTok{se =} \ConstantTok{FALSE}\NormalTok{, }
                \AttributeTok{color =} \FunctionTok{hsv}\NormalTok{(}\FloatTok{0.52}\NormalTok{, }\DecValTok{1}\NormalTok{, }\FloatTok{0.7}\NormalTok{))}
\end{Highlighting}
\end{Shaded}

\begin{figure}[H]

\centering{

\pandocbounded{\includegraphics[keepaspectratio]{SpatAggNet_supp_files/figure-pdf/fig-k-mod-1.pdf}}

}

\caption{\label{fig-k-mod}}

\end{figure}%

\section{Computing species abundance
distributions}\label{computing-species-abundance-distributions}

The package \emph{pika} provides functions for fitting common species
abundance distribution models. I use the \texttt{fitSAD} from
\emph{pika} to preform model selection on three standard SAD forms---the
log-series, Poisson log-normal, and zero-truncated negative
binomial---and record the best fit model and its parameter(s) for each
data set.

\begin{Shaded}
\begin{Highlighting}[]
\CommentTok{\# models to consider}
\NormalTok{mods }\OtherTok{\textless{}{-}} \FunctionTok{c}\NormalTok{(}\StringTok{\textquotesingle{}lseries\textquotesingle{}}\NormalTok{, }\StringTok{\textquotesingle{}plnorm\textquotesingle{}}\NormalTok{, }\StringTok{\textquotesingle{}tnegb\textquotesingle{}}\NormalTok{)}

\CommentTok{\# loop over studies, fitting SAD models}
\NormalTok{sad\_stats }\OtherTok{\textless{}{-}} \FunctionTok{mclapply}\NormalTok{(obs\_dat, }
                      \AttributeTok{mc.cores =}\NormalTok{ nthrd, }
                      \AttributeTok{FUN =} \ControlFlowTok{function}\NormalTok{(x) \{}
\NormalTok{    nsite }\OtherTok{\textless{}{-}} \FunctionTok{nrow}\NormalTok{(x)}
\NormalTok{    x }\OtherTok{\textless{}{-}} \FunctionTok{colSums}\NormalTok{(x)}
\NormalTok{    s }\OtherTok{\textless{}{-}} \FunctionTok{fitSAD}\NormalTok{(x, mods)}
    
    \CommentTok{\# which model is best supported}
\NormalTok{    i }\OtherTok{\textless{}{-}} \FunctionTok{which.min}\NormalTok{(}\FunctionTok{sapply}\NormalTok{(s, AIC))}
    
    \CommentTok{\# mle of params of bets model}
\NormalTok{    o }\OtherTok{\textless{}{-}}\NormalTok{ s[[i]]}\SpecialCharTok{$}\NormalTok{MLE}
    
    \CommentTok{\# if lseries is best, only has 1 param}
    \ControlFlowTok{if}\NormalTok{(i }\SpecialCharTok{==} \DecValTok{1}\NormalTok{) o }\OtherTok{\textless{}{-}} \FunctionTok{c}\NormalTok{(o, }\ConstantTok{NA}\NormalTok{)}
    
\NormalTok{    o }\OtherTok{\textless{}{-}} \FunctionTok{c}\NormalTok{(i, o, }\FunctionTok{sum}\NormalTok{(x), }\FunctionTok{length}\NormalTok{(x), nsite)}
    \FunctionTok{names}\NormalTok{(o) }\OtherTok{\textless{}{-}} \ConstantTok{NULL}
    
    \FunctionTok{return}\NormalTok{(o)}
\NormalTok{\}) }\SpecialCharTok{|\textgreater{}} 
    \FunctionTok{do.call}\NormalTok{(rbind, }\AttributeTok{args =}\NormalTok{ \_) }\SpecialCharTok{|\textgreater{}} 
    \FunctionTok{as.data.frame}\NormalTok{() }

\NormalTok{sad\_stats }\OtherTok{\textless{}{-}} \FunctionTok{cbind}\NormalTok{(}\AttributeTok{study =} \FunctionTok{rownames}\NormalTok{(sad\_stats), sad\_stats)}
\FunctionTok{names}\NormalTok{(sad\_stats)[}\SpecialCharTok{{-}}\DecValTok{1}\NormalTok{] }\OtherTok{\textless{}{-}} \FunctionTok{c}\NormalTok{(}\StringTok{\textquotesingle{}mod\textquotesingle{}}\NormalTok{, }\StringTok{\textquotesingle{}par1\textquotesingle{}}\NormalTok{, }\StringTok{\textquotesingle{}par2\textquotesingle{}}\NormalTok{, }\StringTok{\textquotesingle{}J\textquotesingle{}}\NormalTok{, }\StringTok{\textquotesingle{}nspp\textquotesingle{}}\NormalTok{, }\StringTok{\textquotesingle{}nsite\textquotesingle{}}\NormalTok{)}
\FunctionTok{rownames}\NormalTok{(sad\_stats) }\OtherTok{\textless{}{-}} \ConstantTok{NULL}
\NormalTok{sad\_stats}\SpecialCharTok{$}\NormalTok{mod }\OtherTok{\textless{}{-}}\NormalTok{ mods[sad\_stats}\SpecialCharTok{$}\NormalTok{mod]}
\end{Highlighting}
\end{Shaded}

\begin{Shaded}
\begin{Highlighting}[]
\NormalTok{sad\_select }\OtherTok{\textless{}{-}} \FunctionTok{filter}\NormalTok{(sad\_stats) }\SpecialCharTok{|\textgreater{}} 
    \FunctionTok{mutate}\NormalTok{(}\AttributeTok{abund\_bin =} \FunctionTok{cut}\NormalTok{(J, }\FunctionTok{c}\NormalTok{(}\FunctionTok{seq}\NormalTok{(}\DecValTok{1}\NormalTok{, }\DecValTok{1400}\NormalTok{, }\AttributeTok{length.out =} \DecValTok{9}\NormalTok{), }
                                \FunctionTok{exp}\NormalTok{(}\FunctionTok{seq}\NormalTok{(}\FunctionTok{log}\NormalTok{(}\DecValTok{2000}\NormalTok{), }\FunctionTok{log}\NormalTok{(}\FunctionTok{max}\NormalTok{(J)), }
                                        \AttributeTok{length.out =} \DecValTok{4}\NormalTok{))))) }\SpecialCharTok{|\textgreater{}} 
    \FunctionTok{group\_by}\NormalTok{(abund\_bin) }\SpecialCharTok{|\textgreater{}} 
    \FunctionTok{summarize}\NormalTok{(}\AttributeTok{lseries\_max =} \FunctionTok{mean}\NormalTok{(mod }\SpecialCharTok{==} \StringTok{"lseries"}\NormalTok{),}
              \AttributeTok{plnorm\_max =} \FunctionTok{mean}\NormalTok{(mod }\SpecialCharTok{==} \StringTok{"plnorm"}\NormalTok{) }\SpecialCharTok{+}\NormalTok{ lseries\_max, }
              \AttributeTok{tnegb\_max =} \FunctionTok{mean}\NormalTok{(mod }\SpecialCharTok{==} \StringTok{"tnegb"}\NormalTok{) }\SpecialCharTok{+}\NormalTok{ plnorm\_max, }
              \AttributeTok{lseries\_min =} \DecValTok{0}\NormalTok{, }
              \AttributeTok{plnorm\_min =}\NormalTok{ lseries\_max,}
              \AttributeTok{tnegb\_min =}\NormalTok{ plnorm\_max) }\SpecialCharTok{|\textgreater{}} 
    \FunctionTok{ungroup}\NormalTok{() }\SpecialCharTok{|\textgreater{}} 
    \FunctionTok{pivot\_longer}\NormalTok{(}\AttributeTok{cols =} \SpecialCharTok{{-}}\NormalTok{abund\_bin,}
                 \AttributeTok{names\_to =} \FunctionTok{c}\NormalTok{(}\StringTok{"model"}\NormalTok{, }\StringTok{".value"}\NormalTok{),}
                 \AttributeTok{names\_sep =} \StringTok{"\_"}\NormalTok{) }\SpecialCharTok{|\textgreater{}} 
    \FunctionTok{mutate}\NormalTok{(}\AttributeTok{abund\_max =} \FunctionTok{gsub}\NormalTok{(}\StringTok{".*,|}\SpecialCharTok{\textbackslash{}\textbackslash{}}\StringTok{]"}\NormalTok{, }\StringTok{""}\NormalTok{, abund\_bin) }\SpecialCharTok{|\textgreater{}}
               \FunctionTok{as.numeric}\NormalTok{(),}
           \AttributeTok{abund\_min =} \FunctionTok{gsub}\NormalTok{(}\StringTok{"}\SpecialCharTok{\textbackslash{}\textbackslash{}}\StringTok{(|,.*"}\NormalTok{, }\StringTok{""}\NormalTok{, abund\_bin) }\SpecialCharTok{|\textgreater{}}
               \FunctionTok{as.numeric}\NormalTok{())}

\FunctionTok{ggplot}\NormalTok{(sad\_select, }\FunctionTok{aes}\NormalTok{(}\AttributeTok{xmin =}\NormalTok{ abund\_min, }\AttributeTok{xmax =}\NormalTok{ abund\_max, }
                       \AttributeTok{ymin =}\NormalTok{ min, }\AttributeTok{ymax =}\NormalTok{ max, }\AttributeTok{fill =}\NormalTok{ model)) }\SpecialCharTok{+}
    \FunctionTok{geom\_rect}\NormalTok{(}\AttributeTok{color =} \StringTok{"black"}\NormalTok{, }\AttributeTok{linewidth =} \FloatTok{0.5}\NormalTok{) }\SpecialCharTok{+}
    \FunctionTok{scale\_fill\_discrete}\NormalTok{(}\AttributeTok{name =} \StringTok{"Model favored}\SpecialCharTok{\textbackslash{}n}\StringTok{by AIC"}\NormalTok{, }
                        \AttributeTok{labels =} \FunctionTok{c}\NormalTok{(}\StringTok{"Log series"}\NormalTok{, }
                                   \StringTok{"Poisson log normal"}\NormalTok{, }
                                   \StringTok{"Truncated negative binomial"}\NormalTok{),}
                        \AttributeTok{palette =} \FunctionTok{hsv}\NormalTok{(}\FunctionTok{c}\NormalTok{(}\FloatTok{0.8}\NormalTok{, }\FloatTok{0.7}\NormalTok{, }\FloatTok{0.52}\NormalTok{), }
                                      \FunctionTok{c}\NormalTok{(}\FloatTok{0.7}\NormalTok{, }\FloatTok{0.8}\NormalTok{, }\DecValTok{1}\NormalTok{), }
                                      \FunctionTok{c}\NormalTok{(}\FloatTok{0.5}\NormalTok{, }\FloatTok{0.9}\NormalTok{, }\FloatTok{0.7}\NormalTok{))) }\SpecialCharTok{+}
    \FunctionTok{xlab}\NormalTok{(}\StringTok{"Total individuals"}\NormalTok{) }\SpecialCharTok{+}
    \FunctionTok{ylab}\NormalTok{(}\StringTok{"Proportion"}\NormalTok{) }\SpecialCharTok{+}
    \FunctionTok{theme\_cowplot}\NormalTok{() }\SpecialCharTok{+}
    \FunctionTok{theme}\NormalTok{(}\AttributeTok{legend.position =} \StringTok{"top"}\NormalTok{) }\SpecialCharTok{+}
    \FunctionTok{guides}\NormalTok{(}\AttributeTok{fill =} \FunctionTok{guide\_legend}\NormalTok{(}\AttributeTok{nrow =} \DecValTok{3}\NormalTok{))}
\end{Highlighting}
\end{Shaded}

\begin{figure}[H]

\centering{

\pandocbounded{\includegraphics[keepaspectratio]{SpatAggNet_supp_files/figure-pdf/fig-aic-sad-1.pdf}}

}

\caption{\label{fig-aic-sad}Proportion of studies favoring one of three
possible species abundance distribution models: log series, Poisson log
normal, truncated negative binomial. Model preference is shown across
different different total numbers of individuals recorded by each study.
Bins start at a width of 175 from 1 to 1,400, then expand in width
exponentially. We see there is no consistent pattern for which model is
favored across total number of individuals.}

\end{figure}%

\section{Studying networks simulated based on real abundance
data}\label{studying-networks-simulated-based-on-real-abundance-data}

Now we seek to compare the networks reconstructed from real data to
those reconstructed from purely random data simulated by sampling SADs
and SSADs without any true interaction/association between species.

First we compute network statistics for the real data

\begin{Shaded}
\begin{Highlighting}[]
\NormalTok{obs\_nets }\OtherTok{\textless{}{-}} \FunctionTok{mclapply}\NormalTok{(obs\_dat, net\_stats, }\AttributeTok{mc.cores =}\NormalTok{ nthrd) }\SpecialCharTok{|\textgreater{}} 
    \FunctionTok{do.call}\NormalTok{(rbind, }\AttributeTok{args =}\NormalTok{ \_) }\SpecialCharTok{|\textgreater{}} 
    \FunctionTok{as.data.frame}\NormalTok{() }\SpecialCharTok{|\textgreater{}} 
\NormalTok{    (\textbackslash{}(x) \{}
\NormalTok{        x }\OtherTok{\textless{}{-}} \FunctionTok{cbind}\NormalTok{(}\AttributeTok{study =} \FunctionTok{rownames}\NormalTok{(x), x)}
        \FunctionTok{rownames}\NormalTok{(x) }\OtherTok{\textless{}{-}} \ConstantTok{NULL}
\NormalTok{        x}
\NormalTok{    \})()}
\end{Highlighting}
\end{Shaded}

\begin{Shaded}
\begin{Highlighting}[]
\FunctionTok{ggplot}\NormalTok{(}\FunctionTok{filter}\NormalTok{(obs\_nets, }\SpecialCharTok{!}\FunctionTok{is.na}\NormalTok{(pos\_rho) }\SpecialCharTok{\&} \SpecialCharTok{!}\FunctionTok{is.na}\NormalTok{(neg\_rho)), }
       \FunctionTok{aes}\NormalTok{(pos\_wm\_abund }\SpecialCharTok{{-}}\NormalTok{ neg\_wm\_abund)) }\SpecialCharTok{+}
    \FunctionTok{geom\_histogram}\NormalTok{(}\AttributeTok{bins =} \DecValTok{20}\NormalTok{, }\AttributeTok{alpha =} \FloatTok{0.5}\NormalTok{, }\AttributeTok{color =} \StringTok{"black"}\NormalTok{) }\SpecialCharTok{+}
    \FunctionTok{xlim}\NormalTok{(}\SpecialCharTok{{-}}\FloatTok{0.33}\NormalTok{, }\FloatTok{0.33}\NormalTok{)}


\FunctionTok{ggplot}\NormalTok{(}\FunctionTok{filter}\NormalTok{(obs\_nets, }\SpecialCharTok{!}\FunctionTok{is.na}\NormalTok{(pos\_rho) }\SpecialCharTok{\&} \SpecialCharTok{!}\FunctionTok{is.na}\NormalTok{(neg\_rho)), }
       \FunctionTok{aes}\NormalTok{(pos\_rho }\SpecialCharTok{{-}}\NormalTok{ neg\_rho)) }\SpecialCharTok{+}
    \FunctionTok{geom\_histogram}\NormalTok{(}\AttributeTok{bins =} \DecValTok{20}\NormalTok{, }\AttributeTok{alpha =} \FloatTok{0.5}\NormalTok{, }\AttributeTok{color =} \StringTok{"black"}\NormalTok{) }\SpecialCharTok{+}
    \FunctionTok{xlim}\NormalTok{(}\FunctionTok{c}\NormalTok{(}\SpecialCharTok{{-}}\DecValTok{2}\NormalTok{, }\DecValTok{2}\NormalTok{))}

\FunctionTok{ggplot}\NormalTok{(}\FunctionTok{filter}\NormalTok{(obs\_nets, }\SpecialCharTok{!}\FunctionTok{is.na}\NormalTok{(pos\_rho) }\SpecialCharTok{\&} \SpecialCharTok{!}\FunctionTok{is.na}\NormalTok{(neg\_rho)), }
       \FunctionTok{aes}\NormalTok{(neg\_rho)) }\SpecialCharTok{+}
    \FunctionTok{geom\_histogram}\NormalTok{(}\AttributeTok{bins =} \DecValTok{24}\NormalTok{, }\AttributeTok{alpha =} \FloatTok{0.5}\NormalTok{, }\AttributeTok{color =} \StringTok{"black"}\NormalTok{) }\SpecialCharTok{+}
    \FunctionTok{xlim}\NormalTok{(}\SpecialCharTok{{-}}\DecValTok{1}\NormalTok{, }\DecValTok{1}\NormalTok{)}
\end{Highlighting}
\end{Shaded}

\begin{Shaded}
\begin{Highlighting}[]
\NormalTok{sim\_nets\_dat }\OtherTok{\textless{}{-}} \FunctionTok{sim\_nets}\NormalTok{(sad\_stats, }
                         \AttributeTok{mc\_cores =}\NormalTok{ nthrd, }
                         \AttributeTok{k\_fun =}\NormalTok{ k\_fun, }
                         \AttributeTok{nsim =} \FunctionTok{round}\NormalTok{(}\FloatTok{1.25} \SpecialCharTok{*} \FunctionTok{sum}\NormalTok{(}\SpecialCharTok{!}\FunctionTok{is.na}\NormalTok{(obs\_nets}\SpecialCharTok{$}\NormalTok{pos\_rho))))}
\end{Highlighting}
\end{Shaded}

\section*{References}\label{references}
\addcontentsline{toc}{section}{References}

\phantomsection\label{refs}
\begin{CSLReferences}{1}{1}
\bibitem[\citeproctext]{ref-quarto}
Allaire, J.J., Teague, C., Scheidegger, C., Xie, Y., Dervieux, C. \&
Woodhull, G. (2025).
\href{https://doi.org/10.5281/zenodo.5960048}{{Quarto}}.

\bibitem[\citeproctext]{ref-calatayud2019}
Calatayud, J., Andivia, E., Escudero, A., Melián, C.J., Bernardo-Madrid,
R., Stoffel, M., Aponte, C., Medina, N.G., Molina-Venegas, R., Arnan,
X., Rosvall, M., Neuman, M., Noriega, J.A., Alves-Martins, F., Draper,
I., Luzuriaga, A., Ballesteros-Cánovas, J.A., Morales-Molino, C.,
Ferrandis, P., Herrero, A., Pataro, L., Juen, L., Cea, A. \&
Madrigal-González, J. (2019). Positive associations among rare species
and their persistence in ecological assemblages. \emph{Nat Ecol Evol}.

\bibitem[\citeproctext]{ref-rcore}
R Core Team. (2018). \emph{\href{https://www.R-project.org/}{R: A
language and environment for statistical computing}}. R Foundation for
Statistical Computing, Vienna, Austria.

\bibitem[\citeproctext]{ref-testthat}
Wickham, H. (2011).
\href{https://journal.r-project.org/archive/2011-1/RJournal_2011-1_Wickham.pdf}{{\emph{testthat}}:
Get started with testing}. \emph{The R Journal}, \textbf{3}, 5--10.

\bibitem[\citeproctext]{ref-roxygen2}
Wickham, H., Danenberg, P. \& Eugster, M. (2018).
\emph{\href{https://CRAN.R-project.org/package=roxygen2}{{\emph{roxygen2}}:
In-line documentation for {R}}}.

\end{CSLReferences}




\end{document}
